\documentclass[../main.tex]{subfiles}

\begin{document}

\subsection{Role of Visualization in the Problem}

The Vietnam Financial Inclusion Dashboard 2025 is designed as an institutional,
five-page interactive policy workspace. Its structure converts 1,000
high-dimensional, weight-adjusted individual records from
\texttt{Findex\_Microdata\_2025\_updateViet Nam.csv} into five connected
thematic layers. That design allows the reader to isolate structural
inequalities, trace transaction flows, and diagnose exclusion barriers across
the Vietnamese financial ecosystem.

\dashboardfigure{0.97}{fig-page1-overview.png}
{Page 1 overview: the macro-socioeconomic baseline.}
{fig:dashboard-overview}

\figref{fig:dashboard-overview}{Page 1 macro-socioeconomic baseline}. The
dashboard begins with a demographic and structural baseline. This page presents
country-level KPI cards such as urban concentration, secondary education
thresholds, and workforce participation, together with cross-tabulated
population views. It appears first because later financial indicators only make
sense when read against the country's underlying social composition.

\dashboardfigure{0.97}{fig-page2-connectivity-overview.png}
{Page 2 overview: connectivity and digital readiness.}
{fig:page2-connectivity-overview}

\figref{fig:page2-connectivity-overview}{Page 2 connectivity overview}. After
the demographic baseline is established, the dashboard moves to connectivity
and digital readiness. This page tracks internet-use depth, smartphone
ownership, and foundational versus online identity coverage so the reader can
determine whether weak financial participation begins as an infrastructure
problem or persists despite high digital access.

\dashboardfigure{0.82}{fig-page3-account-landscape-overview.png}
{Page 3 overview: the default account-ownership and access state.}
{fig:page3-interaction-states}

\figref{fig:page3-interaction-states}{Page 3 default account-access state}. The
narrative then moves to the account ownership and access landscape. This page
reports formal banking and mobile money inclusion, shows how rural and urban
groups distribute across banking tiers, and benchmarks the concrete frictions
faced by the unbanked. It marks the transition from technical readiness to
formal institutional inclusion.

\dashboardfigure{0.97}{fig-page4-cash-flows-overview.png}
{Page 4 overview: cash flows and digitization vectors.}
{fig:page4-cash-flow-overview}

\figref{fig:page4-cash-flow-overview}{Page 4 cash-flow overview}. The next
step tests whether account access actually changes daily transaction behavior.
This page compares cash versus digital delivery across wages, agricultural
sales, transfers, utility payments, and merchant transactions. It helps the
reader move beyond access rates and ask whether account holders remain trapped
in cash-heavy routines.

\dashboardfigure{0.92}{fig-page5-resilience-overview.png}
{Page 5 overview: financial health and resilience.}
{fig:page5-resilience-overview}

\figref{fig:page5-resilience-overview}{Page 5 resilience overview}. The final
page brings the earlier layers together by asking how households cope with
economic shocks. It connects demographics, technology, access, and transaction
behavior to measures of saving, borrowing, and emergency coping channels. This
page is placed last because financial resilience can only be interpreted after
the structural and behavioral context is already clear.

\subsection{Types of Visualization Used}
\label{sec:dashboard-overview}

The dashboard uses a hybrid visualization strategy to match the needs of
policy analysts, development researchers, and decision-makers who need both a
clear narrative and the ability to test subgroup conditions directly.

The explanatory layer functions as an anchored narrative. Its fixed five-page
sequence guides the reader through one policy argument: genuine financial
inclusion cannot be evaluated by account access alone, but must be assessed
through structural context, digital readiness, transaction behavior, and
household resilience. National benchmarks, large labels, and stable page order
keep that storyline legible from beginning to end.

The exploratory layer gives the user room to depart from the national story and
investigate specific demographic conditions. Page-level Power BI filters on
Page 3 and Page 5, together with hover tooltips and linked visual updates, let
the reader ask more precise questions about rural disadvantage, gender gaps, or
income-based exclusion without leaving the broader policy frame.

\begin{table}[H]
\centering
\small
\begin{tabularx}{0.92\textwidth}{p{0.23\textwidth}X}
\toprule
\textbf{Typology} & \textbf{Characteristics} \\
\midrule
Explanatory layer & Fixed five-page narrative; static national benchmarks;
curated policy-facing layout. \\
Exploratory layer & Dynamic multi-criteria slicers; user-led demographic
drill-downs; interactive details-on-demand. \\
\bottomrule
\end{tabularx}
\caption{Hybrid explanatory and exploratory typology used in the dashboard.}
\label{tab:visualization-typology}
\end{table}

\dashboardfigure{0.88}{fig-page3-subgroup-rural-female-poorest20.png}
{Page 3 filtered state for the subgroup Rural + Female + Poorest 20\%.}
{fig:page3-subgroup-state}

Traditional visualization theory often recommends separating exploratory tools
from explanatory presentation so that the user is not pulled away from the main
story too early. This dashboard deliberately relaxes that rule by embedding
multi-criteria filters directly inside the narrative workspace, especially on
\figref{fig:page3-subgroup-state}{Page 3 filtered subgroup state}. That choice
is justified because financial inclusion is inherently intersectional:
policymakers need to isolate cohorts such as rural low-income women and then
observe how the exclusion profile changes relative to the national baseline.
The combination of explanatory sequencing and exploratory filtering therefore
makes the dashboard more useful for real policy diagnosis than either approach
would be on its own.

\end{document}
